% GENERATED FILE. Edit canonical lessons, then run npm run generate:books.
% canonical-source-hash: fnv1a64:9dff982c2f87c0f6
% canonical-lessons: SA-C29-please, SA-C29-pardon, SA-C29-bow, SA-C29-favour, SA-C29-fortunate, SA-R29-second-pass-how-this-book-argues

\chapter{Courtesy Words}
\label{ch:sa-courtesy}
\hlchaptermodality{29}

\begin{chapteropening}
\textbf{By the end of this chapter:} \emph{I can be polite in Sanskrit --- please, pardon, a bow, a favour --- and say what is inside the thank-you I learned on the first day.}

Complete SA-C29-fortunate and demonstrate the chapter promise.
\end{chapteropening}

\section[kṛpayā]{\sk{कृपया} (kṛpayā) — please}

\label{lesson:SA-C29-please}



\begin{quote}
Before the new word: what did \emph{svasti} mean, and what were its two pieces?
\end{quote}

\subsection*{You'll want to know: \sk{कृपया}}
\textbf{\sk{कृपया}} (\emph{kṛpayā}) — \textquotedblleft{}please\textquotedblright{}.

\textbf{\sk{कृपया}} is \emph{please}, and like most polite words in most languages it is not really a word for \textquotedblleft{}please\textquotedblright{} at all. It is the noun \emph{kṛpā}, \textquotedblleft{}compassion, kindness\textquotedblright{}, put into the case that means \emph{by means of}. You are saying \textbf{by your kindness}.

English does the same thing and hides it just as well: \emph{please} was once \emph{if it please you}. Two languages, two buried sentences, one small polite noise.

\begin{grammarlens}[title={What you've built}]
A polite word with a whole clause folded inside it.
\end{grammarlens}

\subsection*{Your turn}
\begin{itemize}
\raggedright
  \item \emph{Say it:} \emph{kṛpayā}
  \item \emph{Say it:} it once more, slowly
  \item \emph{Say it:} \emph{kṛpayā}, and say what it literally claims
  \item \emph{Recall:} say \emph{astu}, then read \textbf{\sk{इदानीम्}}
\end{itemize}

\begin{tcolorbox}[breakable,colback=teal!4,colframe=teal!35!black,title={Before you move on}]
What does \sk{कृपया} mean? (\textquotedblleft{}please\textquotedblright{}.) Say it once more.
\end{tcolorbox}
\section[kṣamyatām]{\sk{क्षम्यताम्} (kṣamyatām) — pardon me; let it be forgiven}

\label{lesson:SA-C29-pardon}



\begin{quote}
Before the new word: what did \emph{kṛpayā} literally claim?
\end{quote}

\subsection*{You'll want to know: \sk{क्षम्यताम्}}
\textbf{\sk{क्षम्यताम्}} (\emph{kṣamyatām}) — \textquotedblleft{}pardon me; let it be forgiven\textquotedblright{}.

\textbf{\sk{क्षम्यताम्}} is the apology. It is longer than \emph{sorry} and it is doing more work: the root means \textquotedblleft{}to be patient, to bear with, to forgive\textquotedblright{}, and the form is a command addressed to nobody in particular. \emph{Let it be forgiven.}

Notice that Sanskrit has now given you three of these impersonal commands: \sk{अस्तु}, \textquotedblleft{}let it be\textquotedblright{}, \sk{क्षम्यताम्}, \textquotedblleft{}let it be forgiven\textquotedblright{}, and \sk{स्वस्ति}, \textquotedblleft{}may it be well\textquotedblright{}. A polite language often prefers to leave the person out.

\begin{grammarlens}[title={What you've built}]
An apology that never names who should do the forgiving.
\end{grammarlens}

\subsection*{Your turn}
\begin{itemize}
\raggedright
  \item \emph{Say it:} \emph{kṣamyatām}
  \item \emph{Say it:} it once more, slowly
  \item \emph{Say it:} \emph{astu}, then \emph{kṣamyatām}, and say what both leave out
  \item \emph{Recall:} read \textbf{\sk{साधु}}, then say \emph{paraśvaḥ}
\end{itemize}

\begin{tcolorbox}[breakable,colback=teal!4,colframe=teal!35!black,title={Before you move on}]
What does \sk{क्षम्यताम्} mean? (\textquotedblleft{}pardon me; let it be forgiven\textquotedblright{}.) Say it once more.
\end{tcolorbox}
\section[praṇāmaḥ]{\sk{प्रणामः} (praṇāmaḥ) — a bow, a respectful salutation}

\label{lesson:SA-C29-bow}



\begin{quote}
Before the new word: what did \emph{kṣamyatām} mean?
\end{quote}

\subsection*{You'll want to know: \sk{प्रणामः}}
\textbf{\sk{प्रणामः}} (\emph{praṇāmaḥ}) — \textquotedblleft{}a bow, a respectful salutation\textquotedblright{}.

\textbf{\sk{प्रणामः}} is a bow — the act, named as a thing. It is \emph{pra} + \emph{nam}, \textquotedblleft{}to bend forward\textquotedblright{}.

And you have met that \emph{nam} twice already without being told. \textbf{\sk{नमस्ते}} is \emph{namas} + \emph{te}, \textquotedblleft{}a bow to you\textquotedblright{}. \textbf{\sk{नमस्कारः}} is \emph{namas} + \emph{kāra}, \textquotedblleft{}the making of a bow\textquotedblright{}. Now \sk{प्रणामः}: the same bending root, a third time, with a prefix. Three of the first words in this book turn out to be one gesture described three ways.

\begin{grammarlens}[title={What you've built}]
A root you had used twice, finally named.
\end{grammarlens}

\subsection*{Your turn}
\begin{itemize}
\raggedright
  \item \emph{Say it:} \emph{praṇāmaḥ}
  \item \emph{Say it:} it once more, slowly
  \item \emph{Say it:} \emph{namaste}, \emph{namaskāraḥ}, \emph{praṇāmaḥ} — and name what all three share
  \item \emph{Recall:} say \emph{alam}, then read \textbf{\sk{यात्रा}}
\end{itemize}

\begin{tcolorbox}[breakable,colback=teal!4,colframe=teal!35!black,title={Before you move on}]
What does \sk{प्रणामः} mean? (\textquotedblleft{}a bow, a respectful salutation\textquotedblright{}.) Say it once more.
\end{tcolorbox}
\section[upakāraḥ]{\sk{उपकारः} (upakāraḥ) — a favour, a good turn}

\label{lesson:SA-C29-favour}



\begin{quote}
Before the new word: what did \emph{praṇāmaḥ} mean, and which root was inside it?
\end{quote}

\subsection*{You'll want to know: \sk{उपकारः}}
\textbf{\sk{उपकारः}} (\emph{upakāraḥ}) — \textquotedblleft{}a favour, a good turn\textquotedblright{}.

\textbf{\sk{उपकारः}} is a favour — something done for someone. Take it apart: \emph{upa}, \textquotedblleft{}near, toward\textquotedblright{}, plus \emph{kāra}, \textquotedblleft{}a doing\textquotedblright{}, from the root behind \textbf{\sk{करोमि}}, \textquotedblleft{}I do\textquotedblright{}.

So a favour is \emph{a doing-toward}. And you have already seen \emph{kāra} once, in \sk{नमस्कारः}, \textquotedblleft{}the making of a bow\textquotedblright{} — the same piece, in a word you have said since the first chapter. Prefixes and pieces are how Sanskrit builds, and by now you can watch it happen.

\begin{grammarlens}[title={What you've built}]
A second word built from the doing-root you already had.
\end{grammarlens}

\subsection*{Your turn}
\begin{itemize}
\raggedright
  \item \emph{Say it:} \emph{upakāraḥ}
  \item \emph{Say it:} it once more, slowly
  \item \emph{Say it:} \emph{karomi}, then \emph{namaskāraḥ}, then \emph{upakāraḥ}
  \item \emph{Recall:} read \textbf{\sk{सत्यम्}}, then say \emph{śubham}
\end{itemize}

\begin{tcolorbox}[breakable,colback=teal!4,colframe=teal!35!black,title={Before you move on}]
What does \sk{उपकारः} mean? (\textquotedblleft{}a favour, a good turn\textquotedblright{}.) Say it once more.
\end{tcolorbox}
\section[dhanyaḥ]{\sk{धन्यः} (dhanyaḥ) — fortunate, blessed}

\label{lesson:SA-C29-fortunate}



\begin{quote}
Before the last word of this run: what did \emph{upakāraḥ} mean, and what did \emph{kṛpayā} mean?
\end{quote}

\subsection*{You'll want to know: \sk{धन्यः}}
\textbf{\sk{धन्यः}} (\emph{dhanyaḥ}) — \textquotedblleft{}fortunate, blessed\textquotedblright{}.

\textbf{\sk{धन्यः}} means \emph{fortunate}, \emph{blessed}, \emph{rich in what matters}. It is built on \emph{dhana}, wealth.

And it closes a circle. The very first thank-you in this book was \textbf{\sk{धन्यवादः}} — and that word is \emph{dhanya} + \emph{vāda}, \textbf{a saying of \textquotedblleft{}you are fortunate\textquotedblright{}}. Sanskrit does not thank you by reporting your gift; it thanks you by calling you lucky. You said the word on the first day and could not see inside it. You can now.

\begin{grammarlens}[title={What you've built}]
Five courtesies, and the first word of this book opened up.
\end{grammarlens}

\subsection*{Your turn}
\begin{itemize}
\raggedright
  \item \emph{Say it:} \emph{dhanyaḥ}
  \item \emph{Say it:} it once more, slowly
  \item \emph{Say it:} all five in order, then take \emph{dhanyavādaḥ} apart into its two pieces
  \item \emph{Recall:} say \emph{evam}, then read \textbf{\sk{स्वस्ति}}
\end{itemize}

\begin{tcolorbox}[breakable,colback=teal!4,colframe=teal!35!black,title={Before you move on}]
What does \sk{धन्यः} mean? (\textquotedblleft{}fortunate, blessed\textquotedblright{}.) Say it once more.
\end{tcolorbox}
\section[dhanyaḥ]{\sk{धन्यः} — the ending you can see, and the argument underneath it}

\label{lesson:SA-R29-second-pass-how-this-book-argues}



\begin{quote}
Three of this chapter's words end on the same soft breath — \emph{praṇāmaḥ}, \emph{upakāraḥ}, \emph{dhanyaḥ}. Say what that breath marks. Then say which of the three names the bow you have been making since your first greeting.
\end{quote}

\subsection*{Script: two marks, read inside words you already say}
Cold, with nothing to help you but the shapes.

Read these aloud:

\begin{itemize}
\raggedright
  \item \textbf{\sk{वैद्यः}} — name the double stroke above the top bar, and say what it replaced
  \item \textbf{\sk{इदानीम्}} — name its opening letter, and say why that vowel is full-size here
\end{itemize}

One of the two is a sign that cannot stand alone and \textbf{replaces} the \emph{a} inside a consonant. The other is the full-size letter a vowel wears when it opens a word and has nothing to lean on. Same vowel sound, two jobs, two shapes.

\begin{culture}
Every claim this book has made about where a word came from has used the same move: \textbf{cut the word at its seam, and compare the pieces.}

\noindent
\begin{tabularx}{\linewidth}{@{}>{\raggedright\arraybackslash}X>{\raggedright\arraybackslash}X@{}}
\toprule
\textbf{the claim} & \textbf{the seam it cuts on} \\
\midrule
Sanskrit, Greek and Latin share a source & verb endings, not vocabulary lists \\
the naming words have cousins west of India & \emph{nāma}, \emph{mama}, \emph{tvam}, \emph{kim} \\
the wellbeing words reach east and west alike & the root travels; the finished word does not \\
leave-taking is built from parts you can name & \emph{again} plus \emph{for seeing} \\
\bottomrule
\end{tabularx}

\emph{Hṛdayam} is the cleanest case the book owns, and \emph{praṇāmaḥ} is this chapter's: \emph{pra} plus the same \emph{nam} that bows in the greeting. \textbf{A word you can take apart is a word you can trace.} A word you cannot take apart — like the ear — stays unplaced no matter how ordinary it is.
\end{culture}

\begin{tcolorbox}[breakable,colback=teal!4,colframe=teal!35!black,title={Before you move on}]
What does the visarga on \emph{dhanyaḥ} mark? (\textbf{Masculine singular.}) Is thanking a close friend in Sanskrit warm, or distancing? (\textbf{Distancing.}) What kind of evidence did the 1786 claim rest on? (\textbf{Grammar.})
\end{tcolorbox}
